\chapter{Product and Technology}\label{ch:product}

\section{What the shopper and the retailer see}\label{sec:product-see}

ClosetTwin is a widget on the retailer's product page and a data feed back to the retailer. The shopper opens the widget, takes two guided photographs with a phone, states how loose or fitted he or she likes the garment to be, and receives a recommended size for the garment on the page. The retailer sees the same process from the other side: it supplies the measurements of each garment, and in return it receives a stream of anonymous fit data, such as which sizes are recommended, which are bought and which come back. Figure~\ref{fig:f3} shows the flow from the scan to the retailer's data. The shopper never sees a three-dimensional avatar, because the first release is about choosing a size and not about visualization (Section~\ref{sec:product-viz}).

\begin{figure}
% F3: scan -> body model + garment specs + fit preference -> size -> retailer data flow.
% No caption/label here: the chapter owns them (fig:f3).
\begin{tikzpicture}[node distance=6mm and 9mm]
  \node[ct node ivory,minimum width=30mm] (scan) {Shopper scan\\(phone photos)};
  \node[ct node ivory,minimum width=30mm,below=of scan] (pref) {Fit preference\\(snug / loose)};
  \node[ct node ivory,minimum width=30mm,below=of pref] (spec) {Garment specs\\(from retailer)};
  \node[ct key,minimum width=26mm,right=14mm of pref] (model) {Body model\\+ fit engine};
  \node[ct node pink,minimum width=22mm,right=10mm of model] (size) {Recommended\\size};
  \node[ct node,minimum width=22mm,right=10mm of size] (shop) {Product page\\widget};
  \node[ct node ivory,minimum width=30mm,below=12mm of size] (feed) {Retailer data feed\\(sizes, returns)};
  \draw[ct arrow] (scan.east) -- ++(5mm,0) |- ([yshift=3mm]model.west);
  \draw[ct arrow] (pref.east) -- (model.west);
  \draw[ct arrow] (spec.east) -- ++(5mm,0) |- ([yshift=-3mm]model.west);
  \draw[ct arrow] (model) -- (size);
  \draw[ct arrow] (size) -- (shop);
  \draw[ct arrow] (shop.south) |- (feed.east);
  \draw[ct arrow,dashed] (feed.west) -| (spec.south);
\end{tikzpicture}

\caption{Data flow: scan, garment specifications and fit preference give a size and a retailer data feed}\label{fig:f3}
\end{figure}

\section{From scan to body model}\label{sec:product-scan}

A body model built from scan data is feasible and well studied, and the studies explain why it is needed. Alhassawi et al. analyzed a dataset of 677 participants and found that only 63 of them (9.15\%) kept the same size across bust, waist and hip, that 614 (90.84\%) fell into different sizes on different measurements, and that 35.45\% were not adequately accommodated by the existing sizing scheme \cite{alhassawi2025}. In other words, a single size label fits a minority of bodies on all dimensions at once. Alemany et al. made the same point earlier, and noted that body-scanning technologies at that time had not achieved a compromise between cost, efficacy and accuracy that would allow wide market penetration \cite{alemany2013}. That study is from 2013 and we cite it as history. What has changed is that phone-based capture now exists: 3DLOOK states that a user can complete a scan independently with two guided smartphone photos and no dedicated hardware \cite{threedlook_pricing}. This is a vendor statement, and it does not tell us how accurate such a scan is for apparel. We found no independent source for the accuracy of phone scans in apparel sizing, so we state it as an engineering target that the pilot must measure and not as a result.

\section{Why scan alone is not enough}\label{sec:product-pref}

A body measurement does not choose a size, because the same body can want different sizes in different garments. Peng and Duraffourg ran an experiment with 175 female participants of shapes from XS to 4XL on six garment tops. They report that fit preference has a significant impact on the size choice, and that knowing the three-dimensional shape of customers alone is not enough for an efficient size recommendation tool \cite{peng2023}. Two further findings, which we know only from their abstracts, point the same way. Sembium et al. report that latent-factor models that use personas and return codes improve the area under the curve by 17 to 21\% against baselines, and that in an online A/B test they raised the percentage of ``fit'' transactions by 0.49\% over the control \cite{sembium2017}. Zhang et al. report that thin models dissuade consumers with larger clothing sizes from buying online because of perceived body-size dissimilarity and a heightened perception of fit risk \cite{zhang2024}. The product therefore takes three inputs and not one: the scan, the garment's measurements and the shopper's stated preference. That the combination beats any single input is a hypothesis of this plan, which the pilot tests (Chapter~\ref{ch:gtm}).

\section{The size model and its accuracy limits}\label{sec:product-model}

The best evidence we hold for the model comes from a study on scan data that used the same kind of measurements. Alhassawi et al. report that a support vector machine that works on the primary measurements reaches an accuracy of 89.66\%, outperforming a model that adds further dimensions through principal-component analysis, which reaches 68.97\% \cite{alhassawi2025}. They also report that sizes 32 and 34 received no correct predictions, which they attribute to overlapping features and too few training samples \cite{alhassawi2025}. We take two lessons from this. First, a simple model on a few good measurements can be accurate enough to be worth testing. Second, the accuracy of such a model is lowest at the edges of the size range, so rare sizes need a fallback, for example a prompt to measure again or a note to the shopper that the recommendation is uncertain. We state the 89.66\% as the best evidence in hand and not as ClosetTwin's measured performance, because the study did not use our scans, our garments or our shoppers, and we have no measurement of our own yet.

\section{Integration with retailer platforms}\label{sec:product-integration}

Integration follows a standard pattern. On Shopify, an app is a set of components that the developer declares in the app configuration, and Shopify runs most of them, so that only the backend is the developer's to host \cite{shopify_dev_build}. A ClosetTwin app on that model needs a widget extension for the product page, a backend for the size model and a feed of the retailer's garment data. We have checked this pattern only for Shopify, which is a minor channel in Israel (Chapter~\ref{ch:gtm}); other platforms are not yet checked, and the plan does not assume that integration there is equally cheap. We assume \assumed{\$0.03} as the cost of one scan in our own infrastructure and \assumed{25\%} of orders that include a scan, which together give a cost to serve of \$0.0075 per order (Chapter~\ref{ch:business}).

\section{Privacy by design: scan for fit, not identification}\label{sec:privacy-design}

A body scan is personal data, and the product is designed so that it is not used for identification. Israel's Privacy Protection Law was amended by Amendment~13, published on 14 August 2024, and the amendment takes effect one year after publication, that is, on 14 August 2025 \cite{il_amend13_2024}. Under the law's definitions, a biometric identifier is data of special sensitivity when it is used, or intended to be used, to identify a person or to verify his or her identity in a digital manner \cite{il_amend13_2024,il_ppl_en}. ClosetTwin scans for fit only. The design therefore never uses body data to identify anyone, collects the minimum that the size model needs, keeps the scan separate from any account identity and, as a design rule, discards the images once the measurements are extracted. We offer explicit consent as \emph{our own design choice}: the law defines consent as informed consent, explicit or implied \cite{il_ppl_en}, so we do not claim that the law requires it. Whether body measurements for sizing could count as health data, which is a separate category of special sensitivity \cite{il_ppl_en}, is an open legal question. We treat it as a design risk that the compliance review must answer, not as a conclusion (Chapter~\ref{ch:risks}).

The regulator's own documents point the same way. The Privacy Protection Authority's 2021 draft opinion, published for public hearing, states that collecting excess data beyond what is necessary is in itself an infringement of privacy \cite{ppa_data_minimization2021}, and its 2025 guide to privacy-enhancing technologies says that building them into the design of systems supports compliance and serves as a foundation for building trust with users \cite{ppa_pets2025}. We report this as the Authority's view. It is not evidence that retailers or shoppers pay more for privacy. The EU's General Data Protection Regulation (GDPR) is secondary for this plan; it applies when Israeli sites sell to EU customers and in the Europe stage, and it treats biometric data for the purpose of uniquely identifying a person as a special category \cite{gdpr2016}.

\section{Visualization: an extension, not the core}\label{sec:product-viz}

Image-based try-on is a separate research field that has been surveyed. Chen et al. reviewed 69 research articles on the factors that influence purchasing decisions and the adoption of virtual try-on \cite{chen2024}, Islam et al. state that the technology has the potential to reduce returns and to optimize conversion \cite{islam2024}, and Song et al. survey methods that synthesize a person image dressed in a given garment \cite{song2023}. The word to keep from this is ``potential'': we know these studies only by their abstracts, and none of them shows a measured reduction in returns. The first release of ClosetTwin therefore leaves visualization out. If the pilot succeeds, a try-on image can be added as an extension on the same data, and Google's try-on model, which we discuss in Chapter~\ref{ch:competition}, shows that this layer is the one most exposed to platform competition.

\ctchapterend
